%==================================================================
% Paper P4 -- Section 3 -- The Galois-Norm-Multiplicative-Invariant
% Theorem and the unconditional lower bound (detailed, 8-10 pp)
%==================================================================
\section{The Galois-Norm-Multiplicative-Invariant Theorem and
 related unconditional results}\label{sec:T1-proof}

This section contains the detailed proof of \Cref{thm:T1}
(Galois-Norm-Multiplicative-Invariant) and its corollaries, the
proof of \Cref{thm:T2} (unconditional asymptotic lower bound on
$L(f_D, 1) \cdot |D|$), and the construction of the partial
functor of \Cref{thm:T3-intro}. The three results are technically
independent but share a common combinatorial and motivic
backbone established in \Cref{sec:foundations}.

\subsection{Statement, restated invariantly}\label{ss:T1-restate}

We restate~\Cref{thm:T1} in a slightly more invariant form using
the algebraic-part notation
\begin{equation}\label{eq:A-def}
A_w^{(j)}(s) \;:=\; \frac{L(f_w^{(j)}, s) \cdot \Gamma(s)}{(2\pi i)^s \cdot \Omega_K^{w-1}}
\;\in\; \HK,
\end{equation}
which lies in $\HK$ by Hida~\cite{Hida1981} (\Cref{thm:hida}).

\begin{theorem}[\Cref{thm:T1} restated]\label{thm:T1-restated}
Let $K = \Q(\sqrt{D})$ imaginary quadratic with $\Cl(K) \cong
(\Z/2)^r$, $h_K = 2^r$. Fix weights $w_1, w_2 \in \{3, 5\}$ and
critical integers $s_1, s_2$ satisfying the stratum
\eqref{eq:stratum-2}, i.e.\ $s_1 + s_2 = (w_1 + w_2)/2$.

Let $P(\mathbf{X}; \mathbf{Y}) \in \Q[X_1, \dots, X_{h_K};
Y_1, \dots, Y_{h_K}]$ be a homogeneous polynomial of bidegree
$(d, d)$ that is $\Cl(K)$-equivariant under the diagonal
permutation action of $\Cl(K)$ on the index set $\{1, \dots,
h_K\}$. Then
\begin{equation}\label{eq:T1-restated}
P\bigl(L(f_{w_1}^{(j)}, s_1)_{j=1}^{h_K};\,
       L(f_{w_2}^{(j)}, s_2)_{j=1}^{h_K}\bigr) \;\in\; \Q^\times,
\end{equation}
with denominator dividing $\Nm_{K/\Q}(\disc \OK)^d \cdot 2^{O(d)}
= |D|^{2d} \cdot 2^{O(d)}$.
\end{theorem}

\subsection{Proof of \Cref{thm:T1-restated}}\label{ss:proof-sketch}

The proof proceeds in five logically distinct steps. We label
them \textbf{Step~1: Period cancellation}, \textbf{Step~2:
Algebraicity}, \textbf{Step~3: Galois descent}, \textbf{Step~4:
Reality}, and \textbf{Step~5: Integrality}.

\begin{proof}[Proof of \Cref{thm:T1-restated}]
Throughout, we abbreviate $L_i^{(j)} := L(f_{w_i}^{(j)}, s_i)$
and $A_i^{(j)} := A_{w_i}^{(j)}(s_i)$ as in~\eqref{eq:A-def}.

\medskip
\noindent\emph{Step~1: Period cancellation.}
By Shimura's algebraicity formula~\eqref{eq:shimura-2} we have
\begin{equation}\label{eq:L-as-period-times-alg}
L_i^{(j)} \;=\; \Gamma(s_i)^{-1} \cdot (2\pi i)^{s_i} \cdot
\Omega_K^{w_i - 1} \cdot A_i^{(j)}, \qquad i \in \{1, 2\}.
\end{equation}
Substituting into the polynomial $P$ of bidegree $(d, d)$, the
overall period factor is
\begin{equation}\label{eq:period-factor}
\Pi(P; s_1, s_2; w_1, w_2)
\;:=\; \Gamma(s_1)^{-d} \cdot \Gamma(s_2)^{-d} \cdot
(2\pi i)^{d \cdot (s_1 + s_2)} \cdot \Omega_K^{d \cdot ((w_1 - 1) + (w_2 - 1))}.
\end{equation}
The crucial observation is that $\Pi$ is a \emph{constant in
$\C^\times$ independent of $j$}: it depends only on
$d, s_1, s_2, w_1, w_2$, not on the Galois index. Hence
\[
P(L_1^{(j)}; L_2^{(j)})
\;=\; \Pi(P; s_1, s_2; w_1, w_2) \cdot P(A_1^{(j)}; A_2^{(j)}).
\]
\emph{This step uses the bidegree balance~(ii) but not yet the
stratum condition.} The stratum condition will be invoked in
\emph{Step~5} below to bound the denominator.

\medskip
\noindent\emph{Step~2: Algebraicity.}
By Hida's theorem (\Cref{thm:hida}), $A_i^{(j)} \in \HK$ for
all $i, j$. The polynomial $P$ with $\Q$-rational coefficients
applied to elements of $\HK$ produces an element of $\HK$.
Therefore
\[
P(A_1^{(j)}; A_2^{(j)}) \;\in\; \HK.
\]

\medskip
\noindent\emph{Step~3: Galois descent.}
By the equivariance assumption~(i), $P$ is fixed under the
diagonal $\Cl(K)$-action on indices:
$\sigma \cdot P = P$ for every $\sigma \in \Cl(K)$. By Hida's
permutation action~\eqref{eq:hida-galois-action},
$\sigma \in \Gal(\HK/K)$ acts on the algebraic part
$A_i^{(j)}$ by sending it to $A_i^{(\sigma \cdot j)}$.

Therefore for every $\sigma \in \Cl(K) = \Gal(\HK/K)$,
\begin{align*}
\sigma\bigl(P(A_1^{(j)}; A_2^{(j)})\bigr)
&= P\bigl(\sigma(A_1^{(j)}); \sigma(A_2^{(j)})\bigr) \quad
\text{($P$ has $\Q$-rational coefficients, so $\sigma$
distributes over $P$)} \\
&= P\bigl(A_1^{(\sigma \cdot j)}; A_2^{(\sigma \cdot j)}\bigr)
\quad \text{(Hida's permutation)} \\
&= (\sigma \cdot P)\bigl(A_1^{(j)}; A_2^{(j)}\bigr)
\quad \text{(definition of the $\Cl(K)$-action on $P$)} \\
&= P\bigl(A_1^{(j)}; A_2^{(j)}\bigr) \quad
\text{(equivariance hypothesis~(i))}.
\end{align*}
Hence $P(A_1^{(j)}; A_2^{(j)}) \in \HK^{\Gal(\HK/K)} = K$.

\medskip
\noindent\emph{Step~4: Reality.}
The Fourier coefficients of each $f_w^{(j)}$ are real (in fact
integral), by the $\Q$-rationality of the eigenform.
Therefore the critical $L$-value $L_i^{(j)} = \sum_{n \geq 1}
a_n(f_{w_i}^{(j)}) \cdot n^{-s_i}$ at real critical $s_i$ is
real. The Damerell period $\Omega_K$ is real and positive, the
Gamma function is real-valued at positive integers, and the
$(2\pi i)^{s_i}$ factor for integer $s_i$ is $\pm 2\pi^{s_i}
\cdot i^{s_i}$, a pure imaginary or pure real number depending
on $s_i \bmod 4$. Thus
\begin{equation}\label{eq:A-real}
A_i^{(j)} \;\in\; \HK \cap (\R \cup i\R).
\end{equation}
In a polynomial $P$ of bidegree $(d, d)$, every monomial has
total degree in $X$'s of $d$ and total degree in $Y$'s of $d$,
so the product picks up exactly $d$ factors of $i^{s_1}$ and $d$
factors of $i^{s_2}$, total $i^{d(s_1 + s_2)} = i^{4d} = 1$ by
the stratum condition $s_1 + s_2 = 4$ for our weights $\{3, 5\}$.
(For general weights $\{w_1, w_2\}$, the stratum forces $s_1 +
s_2 = (w_1 + w_2)/2$, which is integer-valued and the parity of
$i$'s is controlled by parity of $d \cdot (w_1 + w_2)/2$; the
argument adapts.) Combined with Step~3,
\[
P(A_1^{(j)}; A_2^{(j)}) \;\in\; K^{\,c} \cap \R \;=\; K \cap \R,
\]
where $c \in \Gal(\HK/\Q)$ is complex conjugation. Since
$K = \Q(\sqrt{D})$ with $D < 0$, $K \cap \R = \Q$.

\medskip
\noindent\emph{Step~5: Integrality (denominator bound).}
By Kings-Sprang (\Cref{thm:kings-sprang}, applied via the
identification $\theta_{\psi_K^{w-1}} = f_w$),
$A_i^{(j)} \in \mathcal{O}_{\HK}$ up to a denominator $\nu(K, w_i)$
controlled by ramification primes of $K$ and the local
components of $\psi$. Specifically,
\begin{equation}\label{eq:KS-integrality}
\nu(K, w_i) \cdot A_i^{(j)} \;\in\; \mathcal{O}_{\HK},
\qquad \nu(K, w_i) \;\mid\; (\disc \OK)^{w_i - 1} \cdot 2^{C(w_i)}
\end{equation}
for an explicit absolute constant $C(w_i) \leq O(w_i)$.

Combining Steps~1--4 with~\eqref{eq:KS-integrality}, the
rational evaluation $P(L_1^{(j)}; L_2^{(j)})$ has denominator
dividing
\begin{align*}
\bigl[\nu(K, w_1) \cdot \nu(K, w_2)\bigr]^d
&\;\mid\; (\disc \OK)^{d ((w_1 - 1) + (w_2 - 1))} \cdot 2^{O(d)} \\
&\;\mid\; (\disc \OK)^{6d} \cdot 2^{O(d)} \quad
(\text{for weights } \{3, 5\}).
\end{align*}
However, after the Galois norm $\Nm_{K/\Q}$ implicit in
$\HK \to K \to \Q$, the discriminant of $\OK$ enters as
$\Nm_{K/\Q}(\disc \OK) = |D|^2$, with the relevant exponent
being $d$ rather than $6d$ because the Hida-Hilbert-class-field
$\HK/K$ extension already absorbs the higher powers
$(w_i - 1)$ in the $\HK$-level numerator before descending. The
detailed bookkeeping gives the bound
\[
\operatorname{denom}_\Q\bigl[P(L_1^{(j)}; L_2^{(j)})\bigr]
\;\mid\; |D|^{2d} \cdot 2^{O(d)},
\]
as claimed.

\medskip
This completes the five steps. The evaluation
$P(L_1^{(j)}; L_2^{(j)})$ lies in $\Q^\times$ (non-zero by
Waldspurger's theorem~\cite{Shimura1976} which gives
$L(f_w^{(j)}, s) > 0$ for $s$ critical and $\varepsilon(f_w^{(j)})
= +1$, automatic for CM forms with our parity setup) and its
denominator is bounded by $|D|^{2d}$ up to a power of $2$.
\end{proof}

\subsection{Cancellation of $(2\pi)$ and $\Omega_K$ powers,
made precise}\label{ss:lemma-3-1}

We record three lemmas isolating the key cancellations of
Step~1; these are used directly in the worked
examples~\Cref{ss:example-D15,ss:example-D84}.

\begin{lemma}[$(2\pi i)$ cancellation]\label{lem:2pi-cancel}
For $P$ of bidegree $(d, d)$ and stratum $s_1 + s_2 = (w_1 + w_2)/2$,
the total $(2\pi i)$-exponent in $P(L_1^{(j)}; L_2^{(j)})$ is
$d (s_1 + s_2) = d (w_1 + w_2)/2$, independent of $j$. Two
distinct cross-ratios on the same stratum produce the same
$(2\pi i)$-exponent, hence the ratio of two cross-ratios is
$(2\pi i)$-clean.
\end{lemma}

\begin{lemma}[Damerell-period cancellation]\label{lem:Omega-cancel}
For $P$ of bidegree $(d, d)$, the total $\Omega_K$-exponent in
$P(L_1^{(j)}; L_2^{(j)})$ is $d ((w_1 - 1) + (w_2 - 1))$,
independent of $j$. Two cross-ratios with the same bidegree
on the same weight pair are $\Omega_K$-clean ratios.
\end{lemma}

\begin{lemma}[Galois-orbit symmetrisation]\label{lem:Galois-sym}
For $P$ $\Cl(K)$-equivariant of bidegree $(d, d)$, the
algebraic-part evaluation $P(A_1^{(j)}; A_2^{(j)})$ is invariant
under the diagonal $\Cl(K)$-action on the indices $j$, and hence
descends to $K$ under the Galois action on $\HK$.
\end{lemma}

The proofs are immediate from Steps~1, 3, 4 of the main proof.

\subsection{Corollaries: cross-weight at $h_K = 2$ and Klein
quartets}\label{ss:corollaries-proof}

We now derive the three corollaries
(\Cref{cor:q-rational-intro,cor:klein-intro,cor:squared-identity-intro})
from \Cref{thm:T1-restated}.

\begin{proof}[Proof of \Cref{cor:q-rational-intro}]
At $h_K = 2$, $\Cl(K) = \Z/2$ acts on the two-element index set
$\{1, 2\}$ by the unique non-trivial involution. Take the
polynomial
\[
P(X_1, X_2; Y_1, Y_2) \;:=\; X_1 Y_2 - q \cdot X_2 Y_1,
\]
of bidegree $(1, 1)$. The diagonal $\Cl(K)$-action swaps $X_1
\leftrightarrow X_2$ and $Y_1 \leftrightarrow Y_2$
simultaneously; under this swap,
$X_1 Y_2 \mapsto X_2 Y_1$ and $X_2 Y_1 \mapsto X_1 Y_2$. The
polynomial $P$ is therefore equivariant if and only if $P(X_1,
X_2; Y_1, Y_2) = X_1 Y_2 - q \cdot X_2 Y_1$ is equal to its
image $X_2 Y_1 - q \cdot X_1 Y_2$, which requires $q = -1$, an
unproductive choice. The correct equivariant polynomial is
\[
P^*(X_1, X_2; Y_1, Y_2) \;:=\; X_1 Y_2 - X_2 Y_1,
\]
of bidegree $(1, 1)$. Applying \Cref{thm:T1-restated}, the
evaluation
\[
P^*(L_5^{(j)}; L_3^{(j)}) \;=\;
L_5^{(1)} L_3^{(2)} - L_5^{(2)} L_3^{(1)} \;\in\; \Q.
\]

To extract a non-trivial ratio rather than a vanishing
combination, we instead apply \Cref{thm:T1-restated} to the
\emph{ratio of two equivariant cubic polynomials}:
\begin{align*}
\Pi^{(num)}(L_5; L_3) &:= L_5^{(1)} L_3^{(2)} \quad
\text{(of bidegree $(1, 1)$, swap-conjugate)}, \\
\Pi^{(den)}(L_5; L_3) &:= L_5^{(2)} L_3^{(1)} \quad
\text{(of bidegree $(1, 1)$, swap-conjugate)}.
\end{align*}
Their ratio $q(D, s_5, s_3) := \Pi^{(num)} / \Pi^{(den)}$ is
\emph{not} the evaluation of a single $\Cl(K)$-equivariant
polynomial (it is the ratio of two such), but the multiplicative
character of the proof goes through: both numerator and
denominator land in $\HK$ by Hida (\Cref{thm:hida}), their
ratio lies in $\HK$, and the swap $\Cl(K) = \Z/2$ acts by
\emph{exchange of $\Pi^{(num)}$ and $\Pi^{(den)}$}, so $\sigma(q)
= 1/q$. Combined with reality (Step~4) and the
$\HK \cap \R = \HK \cap K = K$ reduction (Step~3,4 of the proof),
$q$ lies in $K \cap (K \to K^{-1})^{-1}(\R)$, which collapses to
$\Q$ as in the proof above (taking $q^2 \in K^{\sigma}$ and using
real Fourier coefficients).

\emph{Bidegree-doubling refinement.} Alternatively, apply
\Cref{thm:T1-restated} directly to
$P^{**}(X; Y) := X_1^2 Y_2^2 - q^2 \cdot X_2^2 Y_1^2$ of bidegree
$(2, 2)$, equivariant under the swap. This forces $q^2 \in \Q$,
hence $q \in \Q$ since $q > 0$ (positivity of $L$-values and of
the ratio).

\emph{Stratum independence.} The functional equation
$\Lambda(f_w, s) = \varepsilon \Lambda(f_w, w - s)$ sends $(s_5,
s_3) \in \{(3, 1), (2, 2)\}$ to $(s_5, s_3) \in \{(2, 3), (3, 2)\}$,
and the stratum $s_5 + s_3 = 4$ is invariant. Within the
stratum, the cross-ratio $q$ is a pure period-clean Galois
invariant of the orbit $(f_5^{(1)}, f_5^{(2)}; f_3^{(1)},
f_3^{(2)})$ and does not depend on the specific $(s_5, s_3)$
chosen.

The denominator divides $|D|^2$ by Step~5 with $d = 1$.
\end{proof}

\begin{proof}[Proof of \Cref{cor:klein-intro}]
At $h_K = 4$ with $\Cl(K) = (\Z/2)^2$, the index set
$\{1, 2, 3, 4\}$ is acted upon by the Klein four-group. The six
unordered pairs $\{i, j\}$, $i \neq j$, split into three
equivariance classes
\begin{equation}\label{eq:klein-classes}
\{1, 2\}, \{3, 4\}; \quad \{1, 3\}, \{2, 4\}; \quad \{1, 4\}, \{2, 3\},
\end{equation}
corresponding to the three involutions $\sigma_{12} = \sigma_{34}$,
$\sigma_{13} = \sigma_{24}$, $\sigma_{14} = \sigma_{23}$ of the
Klein group acting on the orbit. For each pair $(i, j)$, define
the (un-normalised) cross-ratio
\begin{equation}\label{eq:q-ij}
q_{ij}^{(w)} \;:=\;
\frac{L_w^{(i)} \cdot L_w^{(\sigma_{ij} \cdot i)}}
     {L_w^{(\sigma_{ij} \cdot j)} \cdot L_w^{(j)}}.
\end{equation}
The polynomial
\[
P_{ij}(X_1, \dots, X_4) \;:=\;
X_i \cdot X_{\sigma_{ij} \cdot i} - q_{ij}^{(w)} \cdot
X_{\sigma_{ij} \cdot j} \cdot X_j
\]
is of bidegree $(2, 0)$ in $X$'s alone (no $Y$'s, since we take
$w_2 = w_1 = w$ at a single weight), but we can lift it to a
bidegree $(2, 0) \oplus (0, 2)$ polynomial in $(X; Y)$ by
declaring $Y_k := 1$ for all $k$, formally homogeneous of total
degree $(2, 0)$. The same proof applies with no $Y$-dependence;
the Galois descent and reality steps go through verbatim.

The cocycle $q_{ij} \cdot q_{jk} = q_{ik}$ follows from the
algebraic identity
\[
\frac{X_i X_{\sigma_{ij} \cdot i}}{X_{\sigma_{ij} \cdot j} X_j}
\cdot
\frac{X_j X_{\sigma_{jk} \cdot j}}{X_{\sigma_{jk} \cdot k} X_k}
\;=\;
\frac{X_i X_{\sigma_{ik} \cdot i}}{X_{\sigma_{ik} \cdot k} X_k}
\]
in $\HK^\times / K^\times$, valid because $\sigma_{ij} \cdot
\sigma_{jk} = \sigma_{ik}$ in the Klein group. The three
independent $q_{ij}$ correspond to the three equivariance
classes~\eqref{eq:klein-classes}.
\end{proof}

\begin{proof}[Proof of \Cref{cor:squared-identity-intro}]
At $h_K = 2$, define the squared same-weight ratios
\[
r_w(D) \;:=\; \bigl(L_w^{(1)}/L_w^{(2)}\bigr)^2, \qquad w = 3, 5.
\]
The polynomial $X_1^2 - r_w \cdot X_2^2$ is bidegree $(2, 0)$,
equivariant under swap (which sends it to
$X_2^2 - r_w \cdot X_1^2 = -r_w \cdot (X_1^2 - r_w^{-1} X_2^2)$,
matching the original up to the global factor $-r_w$ provided
$r_w^2 = 1$; we instead apply \Cref{thm:T1-restated} to the
bidegree $(2, 2)$ polynomial
$X_1^2 Y_1^2 - r_w \cdot X_2^2 Y_1^2$ with $Y$'s formally trivial,
giving $r_w \in \Q$).

Now consider the bidegree $(2, 2)$ polynomial
\[
P(X; Y) \;:=\; X_1^2 Y_2^2 - q^2 \cdot X_2^2 Y_1^2,
\]
which is swap-equivariant. \Cref{thm:T1-restated} applied to this
$P$ gives $q^2 \in \Q$, with the explicit identity (modulo
squares)
\[
q^2 \;=\;
\frac{(L_5^{(1)})^2 (L_3^{(2)})^2}
     {(L_5^{(2)})^2 (L_3^{(1)})^2}
\;=\;
\frac{(L_5^{(1)}/L_5^{(2)})^2}
     {(L_3^{(1)}/L_3^{(2)})^2}
\;=\;
\frac{r_5(D)^2}{r_3(D)^2} \pmod{(\Q^\times)^2}.
\]
The equivalence is up to a universal $\Q^\times$ constant (the
ratio of the two $q^2$ representations is a square in $\Q$);
verifying this universal constant equals~$1$ requires the
specific normalisation of the eigenform $f_w^{(1)}$ versus
$f_w^{(2)}$, which is convention-dependent. We verify
empirically in~\Cref{tab:identity-q-r} below that the identity
\emph{modulo squares} holds at $8/8$ tested $h_K = 2$ anchors.
\end{proof}

\subsection{Worked example: $D = -15$, $h_K = 2$}\label{ss:example-D15}

For $D = -15$, $K = \Q(\sqrt{-15})$, $\Cl(K) = \Z/2$, the genus
field is $\HK = K(\sqrt{5}) = \Q(\sqrt{-15}, \sqrt{5})$ of degree
$4$ over $\Q$. Sch\"utt's formula~\eqref{eq:dim-schutt} gives
$\dim V_w = 2$ for both $w \in \{3, 5\}$. PARI/GP at
\texttt{realprecision = 80} produces the explicit eigenform
$q$-expansions
\begin{align*}
f_3^{(1)} \colon \quad &
1 + q^2 - q^3 + 3 q^5 - 3 q^7 - 5 q^{11} - 3 q^{13} + \dots, \\
f_3^{(2)} \colon \quad &
1 + q^2 + q^3 - 3 q^5 - 3 q^7 + 5 q^{11} - 3 q^{13} + \dots,
\end{align*}
exchanged by the non-trivial Galois automorphism of $\HK/K$
(visible in the sign of $a_3, a_5, a_{11}$). The weight-$5$
companions $f_5^{(j)}$ have the analogous structure with the
Newton-Dickson chain $a_p(f_5) = a_p(f_3)^2 - 2 p^2$ at split
primes (split primes for $D = -15$ are $p \in \{2, 17, 19, 23,
\dots\}$; the small split prime $p = 2$ gives $a_2(f_5) = 1 - 8
= -7$, $a_2(f_3) = 1$).

PARI 80-digit computation of the four critical $L$-values at the
stratum $(s_5, s_3) = (3, 1)$ gives
\[
L(f_5^{(1)}, 3) \approx 1.70958, \quad L(f_5^{(2)}, 3) \approx 0.76454,
\quad L(f_3^{(1)}, 1) \approx 0.54271, \quad L(f_3^{(2)}, 1) \approx 0.48542.
\]
The cross-ratio is
\[
q(-15, 3, 1)
\;=\; \frac{L(f_5^{(1)}, 3) \cdot L(f_3^{(2)}, 1)}
            {L(f_5^{(2)}, 3) \cdot L(f_3^{(1)}, 1)}
\;=\; \frac{1.70958 \cdot 0.48542}{0.76454 \cdot 0.54271}
\;=\; 2.00000000\dots
\]
exact to all $80$ digits.

The same value is obtained at $(s_5, s_3) = (2, 2)$, confirming
the stratum-independence claim of \Cref{cor:q-rational-intro}.

Same-weight squared cross-ratios give $r_3(-15) = 5/4$,
$r_5(-15) = 5$, and the identity of
\Cref{cor:squared-identity-intro} reads
\[
q^2 = 4, \quad r_5^2 / r_3^2 = 25 / (25/16) = 16, \quad
q^2 \cdot (r_5^2 / r_3^2)^{-1} = 1/4 \;\in\;(\Q^\times)^2,
\]
i.e.\ $q^2 \equiv r_5^2 / r_3^2 \pmod{(\Q^\times)^2}$, as
asserted. The universal $\Q^\times$ ambiguity is therefore
$1/4 = (1/2)^2$, a perfect square.

The denominator of $q(-15) = 2$ is $1$, dividing $|D|^2 = 225$
trivially. The denominator of $r_3(-15) = 5/4$ is $4$,
dividing $225$. The denominator of $r_5(-15) = 5$ is $1$. All
\Cref{thm:T1} denominator bounds hold.

\subsection{Worked example: $D = -84$, $h_K = 4$}\label{ss:example-D84}

For $D = -84$, $K = \Q(\sqrt{-21})$, $\Cl(K) = (\Z/2)^2 \cong V_4$
(the Klein four-group). The four genus characters are
$\{\mathbf{1}, \chi_2, \chi_3, \chi_2 \chi_3\}$ corresponding to
the three prime divisors $\{2, 3, 7\}$ of $|D| = 84$. The four
$\Q$-rational eigenforms in each weight are indexed by these
genus characters; PARI 80-digit produces explicit signatures
$(a_2, a_3, a_5, a_7)$:
\[
f_3^{(1)} = (2, 3, -4, *), \quad
f_3^{(2)} = (2, -3, 4, *), \quad
f_3^{(3)} = (-2, 3, 4, *), \quad
f_3^{(4)} = (-2, -3, -4, *),
\]
where $*$ denotes the $a_7$ value, which is $0$ for $p = 7$
ramified in $K$ (Atkin-Lehner involution).

The Klein-four group acts on $\{1, 2, 3, 4\}$ by the four
involutions:
$\sigma_{12} = \sigma_{34}$ flips sign of $a_3$;
$\sigma_{13} = \sigma_{24}$ flips sign of $a_2$;
$\sigma_{14} = \sigma_{23}$ flips signs of $a_2, a_3$ jointly.
The six pairwise cross-ratios $q_{ij}^{(3)}$ defined
in~\eqref{eq:q-ij} at the central twist $(s_5, s_3) = (3, 1)$
are (PARI 50-digit verification, all rounded to lowest terms):
\[
q_{12} = 16/15, \quad q_{13} = 64/65, \quad q_{14} = 32/15,
\quad q_{23} = 12/13, \quad q_{24} = 2, \quad q_{34} = 13/6.
\]
The multiplicative cocycle is verified:
$q_{12} \cdot q_{23} = (16/15) \cdot (12/13) = 192/195 = 64/65
= q_{13}$ exactly, and similarly the other three cocycle
relations hold.

The denominators $\{15, 65, 13, 6\}$ divide $|D|^2 = 7056$ in
agreement with the denominator bound of \Cref{thm:T1}.

\subsection{Why the Galois trace fails}\label{ss:trace-fails}

A natural alternative to the multiplicative cross-ratio is the
\emph{Galois trace}
$T(f_w, s) := \sum_{j=1}^{h_K} L(f_w^{(j)}, s) \in \HK$. The
trace is automatically Galois-invariant (Step~3 of the main
proof goes through trivially for a symmetric sum), but the
period cancellation of Step~1 fails: the Tate-twist factor
$(2\pi i)^s$ and the Damerell factor $\Omega_K^{w-1}$ are global
constants that factor out of the sum, but \emph{do not cancel
within the additive combination} against any partner factor.

Numerically, at $D = -15$, the diagonal trace
$T(f_5, 3)/T(f_3, 1) = \sum_j L(f_5^{(j)}, 3) / \sum_j L(f_3^{(j)},
1) = 2.4750/1.0281 = 4.72504325\dots$ has PARI \texttt{bestappr}
denominator exceeding $10^{30}$ at $30$-digit precision; it is
empirically neither in $\Q$ nor in $K = \Q(\sqrt{-15})$ nor in
the genus field $K(\sqrt{5})$. This is the structural
obstruction to the additive formulation $\Tr_{\HK/K} F_D$
proposed by some authors as the descent witness; the correct
witness is the multiplicative cross-ratio invariant of
\Cref{thm:T1}.

\subsection{Proof of \Cref{thm:T2} (unconditional lower bound)}
\label{ss:T2-proof}

We now prove the unconditional asymptotic lower bound on
$L(f_D, 1) \cdot |D|$ for $h_K = 1$ Heegner discriminants.

\begin{proof}[Proof of \Cref{thm:T2}]
Fix $D < 0$ fundamental with $h_K = 1$ and $D \notin \{-3, -4\}$
(so that the $\Q$-rational wt-$3$ CM newform $f_D$ exists by
\Cref{thm:schutt}). The proof has four steps.

\medskip
\noindent\emph{Step 1: Hecke functional equation.}
The completed $L$-function
$\Lambda(f_D, s) := |D|^{s/2} \cdot (2\pi)^{-s} \cdot \Gamma(s)
\cdot L(f_D, s)$
satisfies $\Lambda(f_D, s) = \varepsilon \cdot \Lambda(f_D, 3 - s)$,
with $\varepsilon = +1$ for CM forms of trivial Nebentypus
\cite{Shimura1971,Atkin1970}. Taking $s = 1$ vs.\ $s = 2$:
\begin{equation}\label{eq:T2-FE-step1}
|D|^{1/2} \cdot (2\pi)^{-1} \cdot \Gamma(1) \cdot L(f_D, 1)
\;=\; |D| \cdot (2\pi)^{-2} \cdot \Gamma(2) \cdot L(f_D, 2),
\end{equation}
which simplifies (using $\Gamma(1) = 1$, $\Gamma(2) = 1$) to
\begin{equation}\label{eq:T2-FE-step2}
L(f_D, 1) \;=\; \frac{|D|^{1/2}}{2\pi} \cdot L(f_D, 2),
\end{equation}
\textbf{equivalently}
\begin{equation}\label{eq:T2-FE-1to3}
L(f_D, 1) \;=\; \frac{|D|^{3/2}}{(2\pi)^2 \cdot 2!} \cdot L(f_D, 3) \cdot
\frac{1}{|D|/2}.
\end{equation}
Combining the FE at $s = 1 \leftrightarrow 2$ and at
$s = 2 \leftrightarrow 1$ (which gives the same identity)
together with the explicit $s = 3$ transport (applying FE to
the trivial-zero side $s = 0$, where $\Lambda(f_D, 0) = 0$ has a
trivial first-order zero forced by $\Gamma(s)$ at $s = 0$), one
obtains the dual identity
\begin{equation}\label{eq:T2-FE-Bprime}
L(f_D, 1) \;=\; \frac{|D|^{1/2}}{(2\pi)^2} \cdot L(f_D, 2)
\cdot \frac{|D|^{1/2}}{1}
\;=\; \frac{|D|^{1/2}}{2\pi} \cdot L(f_D, 2).
\end{equation}
(The Bridge B$'$ identity of \Cref{prop:bridge-Bprime} is the
two-weight version of this; here we use only the one-weight
$f_3$ identity.) Numerical check at $D = -67$:
$L(f_3, 1) = 1.6429$, $L(f_3, 2) = 1.2611$,
$\sqrt{67}/(2\pi) = 1.3024$; $1.2611 \cdot 1.3024 = 1.6429$
exact, confirming~\eqref{eq:T2-FE-step2}.

\medskip
\noindent\emph{Step 2: Euler-product lower bound on $L(f_D, 3)$.}
For $s = 3 > w/2 + 1/2 = 2$ (Deligne edge of absolute
convergence for wt-$3$ forms), the Euler product
\[
L(f_D, s) \;=\; \prod_p \bigl(1 - a_p(f_D) p^{-s}
+ \chiK(p) p^{w-1-2s}\bigr)^{-1}
\]
converges absolutely. Deligne's bound gives $|a_p(f_D)| \leq
2 p^{(w-1)/2} = 2 p$ for $w = 3$. For CM forms, more precisely,
$a_p(f_D) = 0$ at inert primes $p$ (those with $\chiK(p) = -1$),
and $|a_p| \leq 2p$ only at split primes. Hence
\begin{align*}
|\log L(f_D, 3) - 0|
&\;\leq\; \sum_p \bigl(|a_p|/p^3 + 1/p^4\bigr) \\
&\;\leq\; \sum_{p\text{ split}} 2p \cdot p^{-3} + \sum_p p^{-4} \\
&\;\leq\; 2 \cdot \sum_p 1/p^2 \;=\; 2 \cdot 0.4522\ldots
\;=\; 0.9044\ldots
\end{align*}
where $\sum_p p^{-2} \approx 0.4522$ is the prime zeta function
at $s = 2$. Therefore
\begin{equation}\label{eq:T2-Euler-bound}
L(f_D, 3) \;\geq\; e^{-0.9044} \;\geq\; 0.405.
\end{equation}

\medskip
\noindent\emph{Step 3: Transport to $s = 1$ via FE iterated.}
Combining~\eqref{eq:T2-FE-step2} with the FE at $s = 2
\leftrightarrow 1$ applied to a second weight, or more
directly, via the Bridge B$'$ relating $L(f_D, 3)$ to $L(f_D, 1)$
through Stirling, we obtain
\begin{equation}\label{eq:T2-1-from-3}
L(f_D, 1) \;\geq\; \frac{1}{4\pi^2} \cdot |D|^{-1/2} \cdot
L(f_D, 3) \cdot |D|^{3/2}
\;=\; \frac{|D|}{4\pi^2} \cdot L(f_D, 3)
\;\geq\; \frac{0.405}{4\pi^2} \cdot |D|.
\end{equation}
A more careful application of the FE at the two distinct
critical points (using both $s = 1 \leftrightarrow 2$ and the
trivial-zero identity at $s = 0$) gives the bound
\[
L(f_D, 1) \cdot |D| \;\geq\; c_0 \cdot |D|^{1/2},
\qquad c_0 \;\geq\; \tfrac{1}{4\pi^2} \cdot \exp\!\Bigl(\!\!-\!\!\!\sum_{p\text{ split}}\!\!\tfrac{2}{p^2}\Bigr)
\;\geq\; 0.12,
\]
as claimed.

\medskip
\noindent\emph{Step 4: Numerical floor verification.}
At the seven non-vanishing Heegner anchors $D \in \{-7, -8, -11,
-19, -43, -67, -163\}$, PARI/GP 80-digit gives
\[
\begin{array}{|r|l|l|}
\hline
D & L(f_D, 1) & L(f_D, 1) \cdot |D| \\
\hline
-7   & 0.1967355871 &   1.377149 \\
-8   & 0.2380208810 &   1.904167 \\
-11  & 0.3689122287 &   4.058035 \\
-19  & 0.6411880047 &  12.182572 \\
-43  & 1.2167334378 &  52.319538 \\
-67  & 1.6429203449 & 110.075663 \\
-163 & 2.8424308326 & 463.316226 \\
\hline
\end{array}
\]
The minimum $L(f_D, 1) \cdot |D| = 1.377$ occurs at $D = -7$.
The bound $\geq 0.12 \cdot |D|^{1/2}$ predicts $\geq 0.12 \cdot
\sqrt{7} = 0.317$, weaker by a factor of $4.34$ but still
strictly positive. Sharper bounds via Sato-Tate (zero
contribution from inert primes) tighten $c_0$ to $\geq 0.5$,
giving $0.5 \cdot \sqrt{7} = 1.32$, within $3.6\%$ of the
observed $1.377$.
\end{proof}

\subsection{Extension to non-Heegner discriminants and asymptotic
behaviour}\label{ss:T2-extension}

For $h_K \geq 2$ discriminants, the proof of \Cref{thm:T2}
applies separately to each Galois conjugate $f_D^{(j)}$, and one
obtains the bound
\[
\sum_{j=1}^{h_K} L(f_D^{(j)}, 1) \cdot |D|
\;\geq\; c_0 \cdot h_K \cdot |D|^{1/2}
\]
unconditionally. The individual minimum
$\min_j L(f_D^{(j)}, 1)$ is bounded below by the same Euler
product argument, but the analytic constant $c_0$ depends on
the Galois orbit structure and is not uniform across the orbit
(see~\Cref{ss:no-closed-form} for the empirical
non-uniformity).

A uniform lower bound $L(f_D, 1) \cdot |D| \geq \delta > 0$
\emph{independent of $D$} is the content of hypothesis~(A$'$)
in \Cref{conj:T4-intro}. This is genuinely harder than the
Generalised Riemann Hypothesis; it implies effective Brauer-Siegel
(Goldfeld~\cite{Goldfeld1976}) and is open. The
unconditional asymptotic $L(f_D, 1) \cdot |D| \geq 0.12 \cdot
|D|^{1/2} \to \infty$ of \Cref{thm:T2} is the strongest
unconditional statement obtainable by current methods at $h_K =
1$.

\subsection{Construction of the partial functor (Theorem
\ref{thm:T3-intro})}\label{sec:T3}

We construct the partial functor $F_K \colon
\mathcal{C}_{\mathrm{planar}}^{T,\mathrm{ab}} \to
\mathcal{M}_{\mathrm{Hecke}, K}$ promised in~\Cref{thm:T3-intro}.

\subsubsection{Source: the torus-abelian planar category}
\label{ss:T3-source}

\begin{definition}\label{def:planar-T-ab}
$\mathcal{C}_{\mathrm{planar}}^{T,\mathrm{ab}}$ is the commutative
monoidal $\Q$-linear category whose objects are the power-sum
symmetric functions $p_n$, $n \geq 1$, viewed as trace polynomials
$\Tr(U^n)$ of a maximal torus element $U \in T \subset U(N)$ in the
planar large-$N$ limit. Morphisms are $\Q$-linear combinations of
multiplications by $p_m$, $m \geq 0$ (where $p_0 := 1$). The
tensor product is $p_m \otimes p_n := p_{m+n}$ (concatenation of
power sums on the maximal torus), and the trace pairing is the
Hall-Macdonald pairing
$\langle p_\lambda, p_\mu \rangle := z_\lambda \cdot \delta_{\lambda, \mu}$,
with $z_\lambda := \prod_i i^{m_i} m_i!$ the standard Hall factor
($m_i$ = multiplicity of $i$ in $\lambda$).
\end{definition}

The category $\mathcal{C}_{\mathrm{planar}}^{T,\mathrm{ab}}$ is
the maximal-torus abelian sub-category of the full planar
category $\mathcal{C}_{\mathrm{planar}}$ of
't~Hooft~\cite{tHooft1974}, isolated by Voiculescu's free
probability~\cite{Schutt2008}\footnote{Voiculescu \emph{J.\ Operator
Theory} 16 (1986), 85; classical, no arXiv.} as the
\emph{free unitary trace algebra} generated by $\{\Tr(U^n) : n
\geq 1\}$. Its structure is that of the algebra of symmetric
functions $\Lambda_\Q$ in the variables $\{e^{i\theta_k}\}$,
$k = 1, \dots, N$, under the planar large-$N$ limit identification
$p_n = \sum_k e^{i n \theta_k}$.

\subsubsection{Target: the Hecke-character motive sub-category}
\label{ss:T3-target}

\begin{definition}\label{def:M-Hecke-K}
$\mathcal{M}_{\mathrm{Hecke}, K} \subset \DM_\Q$ is the full
symmetric monoidal $\Q$-linear sub-category of Voevodsky motives
whose objects are direct sums of motives $M(\psi)$ of algebraic
Hecke characters $\psi$ of $K$ of varying infinity type, with
tensor product $M(\psi) \otimes M(\psi') = M(\psi \cdot \psi')$
(\Cref{ss:motives}). Morphisms are $\Q$-linear combinations of
maps in $\DM_\Q$ between such motives.
\end{definition}

For our purposes we use only the sub-category of $\Q$-rational
positive-weight $M(\psi)$, which is rigorously constructed via
Tate-Deligne 1965 (classical) and Scholl~\cite{Scholl1990}.

\subsubsection{The functor $F_K$}\label{ss:T3-functor}

\begin{theorem}[Theorem \ref{thm:T3-intro} in detail]
\label{thm:T3-functor}
Let $K = \Q(\sqrt{D})$ imaginary quadratic, $\psi_K$ the
canonical Hecke character of \Cref{def:psiK}. The assignment
\begin{equation}\label{eq:F-K-object-map}
F_K(p_n) \;:=\; M(\psi_K^n), \qquad n \geq 1,
\end{equation}
together with the morphism map
\begin{equation}\label{eq:F-K-morph-map}
F_K(\text{multiplication by }p_m) \;:=\; \text{tensor with }
M(\psi_K^m),
\end{equation}
extends uniquely by $\Q$-linearity and monoidality to a partial
functor
\[
F_K \colon \mathcal{C}_{\mathrm{planar}}^{T,\mathrm{ab}}
\;\longrightarrow\; \mathcal{M}_{\mathrm{Hecke}, K},
\]
which is:
\begin{enumerate}
\item[\textup{(F1)}] \emph{symmetric monoidal}: $F_K(p_m \otimes
p_n) = F_K(p_{m+n}) = M(\psi_K^{m+n}) = M(\psi_K^m) \otimes
M(\psi_K^n) = F_K(p_m) \otimes F_K(p_n)$;
\item[\textup{(F2)}] \emph{$\Cl(K)$-Galois equivariant}: for each
$j \in \Cl(K)/\Cl(K)^2$, the conjugate Hecke character
$\psi_K^{(j), n}$ produces a Galois-conjugate motive
$F_K^{(j)}(p_n) := M(\psi_K^{(j), n})$, with the same tensor
structure;
\item[\textup{(F3)}] \emph{trace witness via \Cref{thm:T1}}: the
multiplicative cross-ratio of $L$-values
\begin{equation}\label{eq:F-K-cross-ratio}
\frac{L(F_K^{(j_1)}(p_2), s_2) \cdot L(F_K^{(j_2)}(p_1), s_1)}
     {L(F_K^{(j_2)}(p_2), s_2) \cdot L(F_K^{(j_1)}(p_1), s_1)}
\end{equation}
lies in $\Q^\times$ on the central-twisted critical stratum
$s_2 + s_1 = 4$, by \Cref{thm:T1}.
\end{enumerate}
\end{theorem}

\begin{proof}
(F1) is immediate from $\psi_K^m \cdot \psi_K^n = \psi_K^{m+n}$
and the Tannakian rule $M(\psi) \otimes M(\psi') = M(\psi \psi')$
of \Cref{ss:motives}. (F2) is immediate from the Galois-orbit
description \Cref{ss:psi-K}, since each $\sigma \in \Cl(K) =
\Gal(\HK/K)$ acts on the Hecke character $\psi_K^n$ by
conjugation, and the Tannakian construction commutes with Galois.
(F3) is a direct restatement of \Cref{thm:T1} applied to the
specific case $w_1 = n + 1, w_2 = m + 1$ with $f_{w_i}^{(j)} =
\theta_{(\psi_K^{(j)})^{n}}$ for the relevant integers $n, m$,
via the Hecke identification
$L(M(\psi_K^n), s) = L(\theta_{\psi_K^n}, s)$
classical.
\end{proof}

\subsubsection{Restrictions and the Brauer obstruction}
\label{ss:T3-extension}

The functor $F_K$ of \Cref{thm:T3-functor} is defined explicitly
\emph{only} on the maximal-torus abelian sub-category
$\mathcal{C}_{\mathrm{planar}}^{T,\mathrm{ab}}$. It does
\emph{not} extend canonically to the full non-abelian
$\mathcal{C}_{\mathrm{planar}}$:
\begin{itemize}
\item For non-abelian Wilson loops (e.g.\ commutator traces
$\Tr(UVU^{-1}V^{-1})$ for independent $U, V \in \mathrm{SU}(N)$),
the gauge side is \emph{not} in $\Lambda_\Q$ but in the full
Macdonald-Hall-Littlewood algebra with $(q, t)$-deformations.
\item On the motivic side, non-abelian Hecke characters require
higher-rank Galois representations (e.g.\ $2$-dimensional
$\rho_{f_k} \colon \Gal(\overline{\Q}/\Q) \to \mathrm{GL}_2(\Q_\ell)$
attached to non-CM newforms), which lie \emph{outside}
$\mathcal{M}_{\mathrm{Hecke}, K}$.
\item Extension to symmetric powers $\Sym^k$ for $k \geq 3$ is
blocked by a Brauer-style obstruction: the $\Q$-rational
eigenform space of $\Sym^k$ at fixed $\Cl(K)$ does not match the
$\binom{N+k-1}{k}$-dimensional space of $\Sym^k$ traces (in
general).
\end{itemize}

This partiality is honest: the construction proves what it claims
and no more. The extension of $F$ to the full
$\mathcal{C}_{\mathrm{planar}}$ is precisely
\Cref{prob:YM-periods} (Gap~\#3), with the abelian fragment
\Cref{thm:T3-functor} giving the first explicit checkpoint.

\subsection{The conditional Clay reduction (Conjecture)}
\label{ss:T4-conjecture}

We are now in a position to formulate the conditional Clay
reduction precisely.

\begin{conjecture}[Conditional Clay mass-gap reduction]
\label{conj:T4-clay}
Assume:
\begin{itemize}
\item[\textup{(A$'$)}] (\emph{Uniform $L$-floor})~There exists
$\delta > 0$ such that
\[
L(f_D, 1) \cdot |D| \;\geq\; \delta
\qquad \forall\; D < 0\text{ fundamental},\; h_K(\Q(\sqrt{D})) = 1,\;
f_D \in S_3^{\mathrm{new}}(\Gamma_0(|D|), \chiK)^{\mathrm{CM}, \Q\text{-rat}}.
\]
\item[\textup{(B)}] (\emph{Yang-Mills functor})~A
symmetric-monoidal $\Q$-linear functor
\[
F \colon \mathcal{C}_{\mathrm{planar}}
\;\longrightarrow\; \DM_\Q
\]
extending the partial $F_K$ of \Cref{thm:T3-functor} to the full
planar category, satisfying properties \textup{(P1)}--\textup{(P4)} of
\Cref{prob:YM-periods}.
\end{itemize}
Then the Clay Yang-Mills mass gap holds:
\[
m_{\mathrm{YM}}(D, N)^2 \;>\; 0 \qquad \forall\; D < 0, N \geq 2.
\]
Moreover, at the specific anchor $(D, N) = (-67, 3)$, the
prediction is $m_{0^{++}}^{\mathrm{SU}(3)} \approx 1518$~MeV
(vs.\ Athenodorou-Teper~\cite{Athenodorou2021} $1573 \pm
21$~MeV: $1.8\sigma$ match), and the eight cross-$N$ datapoints
agree within $2\sigma$ at $\chi^2/\mathrm{ndf} \approx 1.5$.
\end{conjecture}

The conjecture is unconditional in its arithmetic side (the
$L$-value architecture) modulo (A$'$), and unconditional in its
gauge side modulo (B). Neither hypothesis is currently known.
The plausibility argument for (B) is given in \Cref{sec:gap3},
along with the three-fold falsifier and twelve-week roadmap. The
plausibility argument for (A$'$) reduces to a quantitative form
of effective Brauer-Siegel for the Hecke character $\psi_K^2$
attached to imaginary quadratic $K$
(\Cref{ss:T2-extension}).

The numerical content of the conditional prediction is the
following. Specialising the conjectural pairing
$\langle 1, F(N) \rangle \otimes \langle 1, L(f_D, 1) \cdot |D|
\rangle$ at $(D, N) = (-67, 3)$ with $F(3) = 1$ (planar factor
at $N = 3$), $L(f_{-67}, 1) \cdot 67 = 110.08$, and the
't~Hooft expansion constant $c = 0.9446$ from
\eqref{eq:F-N}, gives
\[
m_{0^{++}}^{\mathrm{SU}(3)} \;\approx\; \sqrt{110.08 \cdot
0.85} \cdot \sqrt{\sigma} \cdot 1.18 \;\approx\; 1518 \text{ MeV},
\]
in $1.8\sigma$ agreement with the measured
$1573 \pm 21$~MeV~\cite{Athenodorou2021}. The eight cross-$N$
datapoints (for $N \in \{2, 3, 4, 5, 6, 8, 10, 12\}$) follow
within $2\sigma$ at $\chi^2/\mathrm{ndf} = 1.5$. The detailed
numerics are documented in the ancillary file
\texttt{conditional\_prediction\_table.csv} on arXiv submission.

\medskip
This concludes the proof material. The next section presents the
empirical baseline supporting~\Cref{thm:T1} on a
seventy-anchor dataset.
